\section{任务概述}
\subsection{引言}
本说明书列出文件备份软件的功能、使用条件和验收标准。项目以 Linux 环境中的目录树为处理对象，提供备份、恢复与快照管理。功能需求使用 \Code{FR-1xx/2xx/3xx/4xx} 编号，质量需求使用 \Code{NFR-xx} 编号，用户用例使用 \Code{UC-xx} 编号，形成统一的需求追踪关系。\E{E02}

\subsection{综合描述}
\subsubsection{产品的状况}
Backup Project 是面向 Linux 的备份软件，从本地备份与恢复逐步扩展为共享 C++17 核心、CLI、Qt 6 桌面前端和远端服务的系统。服务端使用 SQLite 管理元数据，客户端与服务端通过项目实现的协议通信。

\subsubsection{产品的功能}
核心功能是备份目录树、恢复指定还原点以及管理仓库；扩展功能包括规则过滤与预览、多种打包和编码选项、Linux 元数据与特殊文件、自动触发、远程账户和文件级增量链。备份结果既包括归档字节，也包括可用于验证和恢复依赖关系的 manifest 与 identity 副文件。

\subsubsection{用户类和特性}
本地用户通过 CLI 或 GUI 操作自己的数据，关注路径选择、备份结果与恢复安全；远程用户经服务器身份验证和账户认证后操作自己的快照，关注数据隔离与传输保护；服务器管理员在服务机本地使用管理、密钥和证书工具，关注实例状态、存储一致性与维护锁。三类角色分别使用对应的入口与权限。计时器与 inotify 属于系统内部触发机制。

\subsection{运行环境}
\subsubsection{基本配置}
硬件平台为 PC，产品目标为 Linux，已记录验证平台为 Ubuntu 24.04 系列。C++17 与 GNU Make 是构建基础；GUI 需 Qt 6 的相应组件，服务端需 SQLite。产品运行平台为 Linux，Windows 与 macOS 原生运行支持不在本版本实现范围内。
\subsubsection{其他配置}
运行资源需求与源文件规模、条目数量和编码选项有关。运行时需同时容纳归档、编码临时流、恢复暂存目录与增量恢复临时目录；空间预算应按数据集和策略现场检查。远程场景需要可达的 TCP 端点、已配置的 pin 或可信根；SSH 隧道模式还依赖本机 SSH 客户端。备份所需读权限、仓库写权限和设备节点恢复权限均由 Linux 提供。

\subsection{系统范围与限制}
系统处理目录树和受管快照，增量依据文件内容摘要判断变化并保存 delta。当前保存的元数据包括 mode、uid/gid、mtime 及链接、设备信息；ACL、xattr、atime、ctime 和 socket 不在支持范围内。磁盘镜像与数据库事务快照属于其他方案。设备节点和属主恢复受 Linux 权限控制。远端服务支持认证会话恢复及上传、下载断点续传。登录会话的 RESUME 恢复认证状态；上传续传通过服务端检查点和数据摘要验证已接收部分，下载续传通过客户端保存的状态继续传输。定时与实时备份依赖持续运行的 GUI 或 CLI 任务执行进程。

\section{功能需求}
\subsection{功能划分与编号需求}
功能按本地数据处理、交互与自动化、远端服务和工程约束划分。表\ref{tab:functional}列出25条功能需求及其规格，作为模块设计与测试用例的共同依据。
\begin{longtable}{L{0.16\linewidth}L{0.75\linewidth}}
\caption{功能需求规格}\label{tab:functional}\\\toprule 编号 & 功能及验收边界\\\midrule\endfirsthead
\toprule 编号 & 功能及验收边界\\\midrule\endhead
FR-101 & 受管目录备份。扫描源目录并写入仓库，由程序命名，失败不伪造可用归档。（UC-01）\\
FR-102 & 受管快照恢复。完整或增量还原点按格式与依赖链恢复，目标不存在或为空。（UC-02）\\
FR-103 & 仓库配置与管理。配置仓库、列举与删除；损坏归档可诊断，列表不代替完整校验。（UC-01/02）\\
FR-104 & 危险拓扑与覆盖保护。归档不得位于源内，非空恢复目标不覆盖，仓库项不得越界或为软链接。（UC-01/02）\\
FR-105 & 元数据。保存类型、mode、uid/gid、mtime、大小、链接/设备信息；属主恢复受权限限制并报告。（UC-01/02）\\
FR-110 & 规则过滤。11字段；规则内AND、规则间OR、exclude优先；排除目录整体剪枝。（UC-03）\\
FR-111 & 只读预览。预览与备份集合语义相同；展示可截断，扫描不因此提前停止或写状态。（UC-03）\\
FR-120 & 打包解包。MyPack v2、USTAR、fast-ustar；后两者同格式不同I/O，legacy产品只读兼容。（UC-01/02）\\
FR-121 & 压缩解压。支持不压缩、Canonical Huffman 和 Zstandard，采用流式编解码；保留旧 LZSS+Huffman 格式及任务兼容。（UC-01/02）\\
FR-122 & 归档加解密。AES-256-CTR-HMAC或DES-CBC-HMAC；DES教学/兼容，恢复先认证。（UC-04）\\
FR-123 & 归档密码通过交互输入、GUI 会话凭据或受权限保护的密码文件提供；明文密码不写入普通命令行参数、任务配置或日志。（UC-04）\\
FR-130 & 桌面GUI。Qt6 QML七页与保留的Widgets入口，共享核心；忙碌和诊断可见。（UC-01/02/03）\\
FR-140 & 定时备份。任务执行进程按周期评估，无变化时跳过，管理执行历史并保留被依赖的快照。（UC-05/06）\\
FR-141 & 实时备份。inotify 递归监听、事件合并、最大等待和溢出重同步；任务执行进程运行期间有效。（UC-07）\\
FR-142 & 模式组合。手动、定时和实时触发均支持完整及增量策略，复用同一组选项检查规则。（UC-01/05/07）\\
FR-201 & 特殊文件。支持软/硬链接、FIFO、字符/块设备；socket必须显式排除或失败，设备恢复需权限。（UC-01/02）\\
FR-301 & 网络快照操作。认证后上传、列举、下载、删除本人快照，支持上传与下载断点续传；区分原始归档与快照材料包。（UC-08/09）\\
FR-302 & 账户与会话。注册、登录、RESUME、节流、注销与隔离；无状态token使用前回查账户。（UC-10/11）\\
FR-303 & 云端元数据。SQLite记录账户与快照关系；管理员只读查询不创建新库，破坏性维护需数据锁。（UC-12）\\
FR-304 & 安全传输。BPSEC1/2保护BPNET1；配置pin或可信根，不做TOFU、不退回明文或未授权降级。（UC-08/09/13）\\
FR-305 & 文件级增量。MyPack、USTAR、Fast USTAR 均支持完整基线、增量及无变化跳过，并支持不加密、AES、DES 组合；恢复前验证父快照和依赖链。（UC-01/02/08/09）\\
FR-401 & 产品单实例。CLI与GUI共用应用锁；换仓库或配置不能绕过；冲突退出3。（UC-01/02）\\
FR-402 & 自动备份组合。定时和实时任务支持完整、增量及加密；加密任务须具备有效凭据，缺少凭据时提示解锁或报错。（UC-05/07）\\
FR-403 & 快照资产完整性。manifest/identity与归档绑定，保留被依赖祖先，删除按后代优先。（UC-01/02/08/09）\\
FR-404 & 崩溃一致性。单文件同步与原子发布；多文件顺序有边界，state损坏不使正确归档失效。（UC-01/02/06）\\
\bottomrule\end{longtable}


\subsection{系统用例}
系统用例按本地操作和远程操作划分。UC-01 数据备份与 UC-02 数据恢复是两个核心用例；文件过滤、口令输入、自动触发、远端传输与账户维护围绕这两个用例展开。仓库配置、列举和删除属于本地快照管理过程；远端快照列举和依赖安全删除属于远端备份与恢复的管理过程。


手动操作可以交互输入口令；自动任务使用已解锁的 GUI 会话凭据或受权限保护的密码文件。任务配置不保存明文口令。远端操作在账户认证通过后进行，管理员在服务端本机执行维护操作。

\subsubsection{UC-01 数据备份}
用例范围为本地受管快照创建；主执行者是本地用户。用户希望得到可恢复副本，后续恢复者关注内容、元数据与快照依赖的一致性。触发条件是用户执行 \Code{backupctl backup} 或在 GUI 发起备份。

前置条件是源目录可访问、仓库配置可解析、过滤规则和算法组合合法、具有足够写入条件；GUI 还需通过写操作互斥检查。操作成功后生成完整快照或增量文件；符合备份规则的文件无变化时，可以跳过本次增量备份并报告原因。

主成功场景如下：
\begin{enumerate}
\item 用户选择源目录、完整或增量策略、打包/压缩/加密选项以及规则；需要时交互输入本次口令。
\item 系统加载配置，校验仓库、源路径、拓扑与组合，按程序规则分配归档名。
\item 共享扫描器与 Filter 枚举实际备份集合；预览与备份使用相同的扫描和过滤规则，保证文件选择结果一致。
\item 完整策略走打包、可选压缩、可选加密和容器封装；增量策略核对可信基线并产生完整基线、delta 或无变化结果。
\item 写入本次临时产物，执行相应同步与发布步骤；归档发布不覆盖既有同名项，增量路径同时维护副文件与身份关系。
\item 返回归档名、种类或跳过原因；GUI 刷新列表并解除操作占用。
\end{enumerate}

异常或替代流程：源目录不可访问、规则非法、源目录与仓库位置冲突、算法组合不支持或归档写入失败时，操作应明确失败。加密备份缺少有效密码时不得自动改为未加密备份。失败操作不得覆盖已有归档。

产品归档名由程序生成，口令不写入普通命令行参数、任务配置或日志。本用例对应 FR-101、103、104、110、120、121、122、123、305，以及 NFR-01、03、07。

\subsubsection{UC-02 数据恢复}
用例范围为从受管完整或增量还原点重建目录树；主执行者是本地用户。前置条件为快照可读取、名称可安全解析、目标不存在或为空、依赖材料可验证；加密归档需正确口令。操作成功后，内容按所选还原点恢复，权限限制和元数据降级如实报告。失败必须返回失败状态，不能把只应用了部分 delta 的树发布为成功结果。

主成功场景如下：
\begin{enumerate}
\item 用户选择快照和目标目录，需要时输入口令；Catalog 检查名称、类型和仓库边界。
\item 系统解析还原点。增量路径核对父绑定、identity、环和深度，得到 base 到 target 的有序链。
\item 按 magic 识别格式；v2 容器检查头、长度和摘要，加密路径先认证再解密，按声明解压和解包。legacy 与裸 USTAR 按各自路径处理。
\item 在暂存区重建目录、文件与支持的特殊条目；完整恢复按 ownership、mode、mtime 收尾，目录从深到浅处理。
\item 增量路径逐跳提取 delta，重建增量恢复临时目录，先按深度应用删除项，再合并新增和修改项；任一步失败均中止整条链。
\item 全部必要步骤成功后发布目标，返回条目数及诊断；GUI 解除忙碌并显示结果。
\end{enumerate}

异常或替代流程：非空目标不覆盖；错误口令、认证失败、损坏格式、越界路径、缺失父节点、混合链或环均拒绝。非 root 无法恢复任意属主时报告 \Code{skipped_ownership}；需要设备节点权限而无法创建时给出明确错误。恢复失败的状态传播和发布前检查见系统设计与测试报告。

增量恢复只在完整依赖链成功后发布目标；文件系统异常及清理规则见系统设计。

\ReportLandscapeFigure{figures/pdf/01_use_case_local.pdf}{本地备份与恢复用例图}{fig:local-usecases}
\ReportLandscapeFigure{figures/pdf/02_use_case_remote.pdf}{远端备份与恢复用例图}{fig:remote-usecases}

\section{外部接口需求}
\subsection{用户界面}
CLI 提供 backup、restore、preview、repository、config、schedule、realtime、remote 与帮助入口。成功、业务失败、用法错误、实例冲突分别采用退出码 0、1、2、3，错误应包含可定位原因。GUI 以首页、备份、自动备份、备份管理、设置、实时备份、远程备份七个页面组织任务；密码掩码显示，危险删除和注销需确认，任务忙碌时给出冲突原因。源路径和目标路径均以实际输入与核心校验结果为准。

界面统一显示归档名、大小、时间、策略与恢复诊断；预览窗口最多展示 300 条并提示截断，扫描仍按实际规则完成。日期与容量格式复用公共实现，不在 QML 重新解释规则。界面交互通过控制器与共享核心连接，输入校验和错误信息沿用同一实现。

\subsection{硬件接口}
软件通过 Linux 文件系统 API 访问磁盘，通过 socket 访问网络，没有独立硬件设备协议。链接、FIFO、字符及块设备的表达和恢复依赖 Linux 语义；\Code{mknod} 和任意属主设置所需权限不能由应用界面绕过。硬盘容量、可用空间与网络可达性不足时需明确失败。

\subsection{软件接口}
C++17 标准库与 POSIX/Linux API 提供文件、进程、锁、inotify 和随机源；Qt 6 负责窗口、QML、信号槽及前端线程桥接；SQLite C API 管理远端元数据。BPNET1 是业务帧协议，BPSEC1/2 是自定义认证加密传输层，与 TLS 采用不同的协议和格式。SSH 模式通过外部 SSH 进程建立端口转发。配置和自动任务状态使用版本化文本文件，快照数据使用二进制格式。

\section{其他非功能性需求}
\subsection{性能需求}
性能目标是控制编码与传输缓冲区的内存占用，避免文件操作阻塞图形界面：编码与传输采用流式接口，服务端固定 worker 池，长度先校验再分配。默认 worker 数为4、单帧负载上限1 MiB、传输分块256 KiB、默认上传上限8 GiB，均为配置与协议参数；吞吐量、最大用户数及恢复时限尚无专项实测数据。
\subsection{安全性需求}
本地安全优先保护既有用户数据，拒绝危险拓扑、覆盖、路径穿越和无法证明归属的清理。远端需先验证服务器身份，再进行账户认证，限制账户范围，拒绝明文或未授权降级。普通配置和日志不保存明文口令，密码文件由权限限制访问；DES 用于教学与旧格式兼容。
\subsection{软件质量属性}
正确性、失败可见性和可恢复性优先于自动清理程度与运行速度。可维护性依赖共享语义、清晰模块职责和可定位诊断；可复现性依赖固定提交、环境与日志。表\ref{tab:nfr}保留11项正式质量编号及可验证判据。
\begin{longtable}{L{0.23\linewidth}L{0.68\linewidth}}
\caption{非功能性需求与验证方法}\label{tab:nfr}\\\toprule 编号及属性 & 要求与验证依据\\\midrule\endfirsthead
\toprule 编号及属性 & 要求与验证依据\\\midrule\endhead
NFR-01 正确性 & 结构、逐文件字节及摘要一致；增量结果与目标manifest相符。 验证：A-TC-01,02,03,04,24,33\\
NFR-02 元数据正确性 & 权限允许时按声明恢复；跳过属主必须有计数及诊断。 验证：A-TC-12,13\\
NFR-03 失败可见 & 失败返回真实状态，不发布部分内容为成功；无法验证身份则留存诊断。 验证：A-TC-34-42\\
NFR-04 崩溃一致与幂等 & 单文件原子发布，无变化不新建，基线损坏保守重建；多个文件更新由业务流程协调。 验证：A-TC-21,22,26\\
NFR-05 易用性 & CLI退出码区分、GUI忙碌与错误明确、危险操作确认、预览截断可见。 验证：A-TC-05,07,11,27\\
NFR-06 可维护性 & 共享规则/模式/路径语义，核心与Qt分离，源码检查与接口可定位。 验证：E08及接口表\\
NFR-07 本地安全 & 拒绝危险路径、覆盖和无归属依据的删除；文件系统限制见设计章节。 验证：A-TC-06,08,39-42\\
NFR-08 远端安全 & 服务器身份、账户隔离、限速、不降级与数据锁均需检验。 验证：A-TC-28-31\\
NFR-09 可测试与复现 & 记录源码版本、环境、输入、输出、退出码及跳过原因。 验证：E05、E07及E15\\
NFR-10 资源边界 & 流式编码与传输，长度先校验，固定worker；全树索引仍随条目数增长。 验证：脚本边界检查；编解码性能记录见E15\\
NFR-11 构建与安装 & Make 目标可分别构建客户端和服务端，发行包可在目标环境安装和启动。验证：E15 发行包与持续集成检查。\\
\bottomrule\end{longtable}

